\section{Experimental Setup}
\label{sec:experiments}

\subsection{Implementation and Computational Environment}
All experiments are implemented in Python 3.14 utilizing \texttt{scikit-learn} (v1.9.1), \texttt{xgboost} (v3.4.1), \texttt{shap} (v0.52.0), and \texttt{scipy} (v1.18.1). Execution was conducted on an x86\_64 Linux environment equipped with an Intel Core i5-3337U processor @ 1.80\,GHz (2 physical cores, 4 threads) and 4\,GB physical RAM running Ubuntu Linux. Random seed 42 is fixed across data partitioning, cross-validation folds, and subsampling for deterministic reproducibility, with multi-seed evaluations (seeds 42, 123, 456, 789, 2024; which retrain models to vary tree subsampling on the fixed test split) validating variance bounds. Hyperparameters are held at fixed baseline configurations (Table~\ref{tab:model_config}) to evaluate architectural efficacy.

\subsection{Evaluation Metrics and Validation Protocols}
To evaluate detection efficacy on the binary classification task ($1 = \text{Phishing}$, $0 = \text{Legitimate}$), we calculate standard performance metrics: $\text{Accuracy} = \frac{TP + TN}{TP + TN + FP + FN}$, $\text{Precision} = \frac{TP}{TP + FP}$, $\text{Recall} = \frac{TP}{TP + FN}$, $\text{F1-Score} = \frac{2 \times \text{Precision} \times \text{Recall}}{\text{Precision} + \text{Recall}}$, $\text{FPR} = \frac{FP}{FP + TN}$, and $\text{FNR} = \frac{FN}{FN + TP}$, alongside continuous probability diagnostics (ROC-AUC and Average Precision PR-AUC).

Generalization stability is evaluated using 5-Fold Stratified Cross-Validation on the training partition ($N_{\text{train}} = 46,916$), recording the mean $\pm$ standard deviation across folds. Crucially, to prevent data leakage across validation folds, zero-variance feature filtering and standard scaling transformations are fitted strictly within each training fold pipeline rather than globally prior to cross-validation. Furthermore, to avoid test-set reuse in feature selection, candidate subset dimensions were evaluated by 5-fold CV on $\mathcal{D}_{\text{train}}$, selecting $k=20$ as the smallest tested subset within about 0.2 percentage points of full performance (CV F1: 91.85\% at $k=5$, 94.52\% at $k=10$, 95.02\% at $k=20$, 95.17\% at $k=30$, and 95.20\% at $k=98$). Because global SHAP rankings were derived across all of $\mathcal{D}_{\text{train}}$, cross-validation scores for the top-$k$ subsets are slightly optimistic due to feature selection on the full training pool; however, held-out test evaluations remain strictly unaffected. Both Table~\ref{tab:model_comparison} and the $k$-sweep utilize identical 5-fold stratified cross-validation splits and preprocessing pipelines; full-model CV F1 records $95.17\% \pm 0.13\%$ in the multi-model benchmark (evaluated with histogram-based tree splitting, \texttt{tree\_method="hist"}) and $95.20\% \pm 0.17\%$ in the standalone $k$-sweep (evaluated under standard exact tree construction), demonstrating cross-validation stability within 0.03\,pp across splitting routines. Final model evaluation is performed on the untouched held-out test partition ($N_{\text{test}} = 11,729$). Statistical significance between classifier predictions is verified using McNemar's test with Edwards continuity correction.

\subsection{Computational Benchmarking Protocol}
Model inference latency is measured in batched vector mode over the entire held-out test matrix ($N=11,729$) using \texttt{predict()} with multi-threading enabled (\texttt{n\_jobs=-1}). To eliminate cold-start caching artifacts, each model executes two warm-up passes followed by 20 repeated measurement iterations. Latency is standardized as milliseconds per 1,000 queries ($\text{ms} / 10^3\,\text{URLs}$), reporting mean $\pm$ standard deviation. In addition, to account for single-threaded runtime differences across linear and tree models under unbatched stream deployment, we benchmark single-query latency ($N=1$, batch size 1) across 500 repeated trials (Decision Tree: 0.22\,ms, Logistic Regression: 0.30\,ms, XGBoost: 0.71\,ms, Linear SVM: 4.01\,ms, Random Forest: 49.38\,ms; compact Top-20 XGBoost: 0.46\,ms median, 0.56\,ms mean). XGBoost's single-query latency (0.71\,ms) is approximately 60$\times$ the 11.6\,$\mu$s required for in-memory lexical feature extraction; we note that this latency reflects Python interpreter and object conversion overhead, and production deployment via native compiled runtimes (e.g., Treelite or C++ ONNX) can reduce single-query latency into microseconds. Storage footprint is measured by serializing fitted model structures using \texttt{joblib} (compression level 3). In-memory string-based feature extraction speed is benchmarked over 10,000 raw URL string parsing operations (extracting all 84 retained lexical properties including character counts, token delimiters, and structural lengths across URL, domain authority, directory path, filename, and query parameter components in pure CPU memory without I/O), yielding an average extraction latency of 11.6\,$\mu$s per URL (11.6\,ms per $10^3$ URLs).

